\documentclass{article}

%Symbols
\usepackage{fdsymbol}
\usepackage{amsmath}

%Header-Footer
\usepackage{fancyhdr}
\cfoot{\thepage}
\rfoot{C\'esar Hern\'andez Cruz}
\lhead{Tarea 1}
\rhead{\textsc{An\'alisis de Algoritmos}}
%\pagenumbering{gobble}

\renewcommand\headrulewidth{1.5pt}
\makeatletter
\def\headrule{{\if@fancyplain\let\headrulewidth\plainheadrulewidth\fi
\hrule\@height\headrulewidth\@width\headwidth
\vskip-\headrulewidth
\vskip-1.5pt}}
\makeatother

%Margins
\usepackage{geometry}
%The cake is a lie
\addtolength{\hoffset}{-0.2cm}
\addtolength{\textwidth}{0.4cm}
\addtolength{\voffset}{-0.5cm}
\addtolength{\textheight}{1cm}

%Algorithms
\usepackage[ruled,vlined,linesnumbered,onelanguage]{algorithm2e}
\SetKwInput{KwIn}{Entrada}
\SetKwInput{KwOut}{Salida}
\SetKw{KwAnd}{\textbf{and}}
\SetKw{KwOr}{\textbf{or}}
\SetKw{KwExchange}{intercambia}


\pagestyle{fancyplain}

%User defined
\newcommand{\Sol}{\textbf{Soluci\'on:} \ }
\newcommand{\set}[1]{\left\{ #1 \right\}}
\newcommand{\setdef}[2]{\left\{ #1 \colon\ #2 \right\}}
\newcommand{\brac}[1]{\left( #1 \right)}


\begin{document}
\begin{enumerate}
    \item Considere la siguiente variante de \textsc{Quicksort}.

        \begin{minipage}[H]{0.943\textwidth}
        \begin{algorithm}[H]
            \DontPrintSemicolon
            \KwIn{Un arreglo $A$ y dos enteros $p$ y $r$.}
            \KwOut{El mismo arreglo $A$, ordenado.}
            \SetKwProg{Fn}{Algoritmo}{:}{}
            \Fn{$\textsc{OtroQuicksort}(A, p, r)$}{
                \While{$p < r$}{
                    $q \gets \textsc{Particion}(A, p, r)$\;
                    $\textsc{OtroQuicksort}(A,p,q-1)$\;
                    $p \gets q+1$
                }
            }
            \caption{\textsc{OtroQuicksort}}
        \end{algorithm}
        \end{minipage}

        \begin{enumerate}
            \item Explique en alto nivel qu\'e hace \textsc{OtroQuicksort} y
                en qu\'e difiere de la versi\'on cl\'asica.

            Este algoritmo es muy similar al Quicksort original, excepto que este en lugar de hacer dos llamadas recursivas
            para ordenar las dos mitades del arreglo (mitades que están dadas por el índice del pivote q al igual que el algoritmo originla), solo hace la
            llamada recursiva para ordenar la mitad de la izquierda y después hace la iteración del ciclo while para ordenar
            segunda mitad (derecha) ya con el índice p actualizado a q+1. Y nuevamete esta segunda mitad es ordenada de la misma manera.
            
            \medskip
            \textbf{**Aclaración**: Cuando digo mitades quiero decir subarreglos, para que no haya confusión.}
            \medskip

            \item Demuestre que al ejecutar $\textsc{OtroQuicksort}(A, 0,
                |A|-1)$ el arreglo $A$ se ordena correctamente.

            Bueno pues puedo mostrar que funciona el algoritmo por inducción siguiendo la longitud del arreglo:

            \textbf{Caso base:} Si el arreglo tiene longitud 1 (o menor), entonces p = r (o en otro caso $p > r$) y el ciclo while no se ejecuta, por lo que el arreglo queda ordenado.
            Y en realidad no impora que el algoritmo no se ejecute en este paso pues cómo sabemos, un arreglo de longitud menor o igual a 1 ya está ordenado.
            
            \medskip
            
            
            \textbf{HI} Supongamos que el algoritmo funciona para arreglos de longitud menor a n.

            \medskip
            \textbf{Paso inductivo:} 
            Ahora tomemos un arreglo A de longitud n. Al incio de ejecutar el alg. se llama a la función Particion(A, p, r) que nos devuelve un índice q tal que todos los elementos
            a la izquierda de q son menores o iguales al elemento en A[q] y todos los elementos a la derecha de q son mayores o iguales al elemento A[q] (justo como lo vimos en clase).
            Ahora el sig. paso es hacer la llamada recursiva a OtroQuicksort(A, p, q-1) que por la hipótesis inductiva sabemos que ordenará correctamente el subarreglo A[p..q-1] ya
            que su longitud es menor a n (ya que incluso en el peor de los casos quitamos solamente al pivote del arreglo).
            Ahora el siguiente paso es actualizar p a q+1 y volver a ejecutar el ciclo while, lo que nos llevará a hacer la llamada recursiva a OtroQuicksort(A, q+1, r) que por la
            hipótesis inductiva sabemos que ordenará correctamente el subarreglo A[q+1..r] ya que su longitud es menor a n (ya que incluso en el peor de los casos quitamos solamente
            al pivote del arreglo). Por lo tanto, al final de ejecutar el algoritmo, tendremos que A[p..q-1] y A[q+1..r] están ordenados y además todos los elementos de A[p..q-1] son 
            menores o iguales al elemento A[q] y todos los elementos de A[q+1..r] son mayores o iguales al elemento A[q] (el cual ya está en su lugar desde la primer ejecución de Particion), por lo que podemos concluir que el arreglo A está ordenado correctamente.

            \medskip
            Entonces en particular se cumple para...$\textsc{OtroQuicksort}(A, 0,
                |A|-1)$

            \item Describa un caso donde la altura de la recursi\'on al ejecutar
                sea $\Theta(n)$ para un arreglo de longitud $n$.  Justifique su
                respuesta.

            Un caso dónde la alura de la pila de recursión sea $\Theta(n)$ es justamente el peor caso de
            Quicksort, querer ordenar un arreglo que ya está ordenado (ya sea de manera ascendente o descendente)
            y que el pivote elegido sea siempre el último o el primer elemento del arreglo.
            
            \medskip
            Ya que cuando acamodamos el pivote en su posición (SPG en la posición final) del lado izquierdo nos
            quedan n-1 elementos, a los cuales les volvemos a aplicar OtroQuicksort (del lado izquierdo nos quedan n-2 elem.) y así sucesivamente hasta que nos quedemos con un subarreglo de longitud 1.
            Lo que nos deja una altura de la pila de recursión de n (contando la llamada incial), que es $\Theta(n)$.

            
            \item Sin alterar el tiempo promedio de ejecuci\'on (es decir, $O(n
                \log n)$), modifique \textsc{OtroQuicksort} para que la altura
                de la recursi\'on sea $\Theta(\log n)$ en el peor caso.

            Como sabemos que OtroQuickSort va apilando llamadas recursivas primero unicamente
            del lado izquierdo, y también sabemos que en el peor de los casos del algoritmo
            se van apilando k-1 llamadas en cada paso k. Y también sabemos que este inciso
            nos pide que la altura de la pila de recursión sea $\Theta(\log n)$ en el peor caso, para
            lograr esto debemos hacer que el agoritmo solamente apile llamadas cuando el tamaño del
            subarreglo sea a lo más la mitad del tamaño de n. Con esto lograremos una "convergencia
            logaritmica" (por así decirlo) de la altura. Ya que si emopezamos a dividir n/2 hasta
            llegar a 1, estos nos llevara justamente $\log n$ pasos. oOk, entonces para lograr esto, podemos
            enviar únicamente el subarreglo pequeño a la recursión pues sabemos que este va a tener a lo más
            n/2 elementos (es el caso dónde ambos subarreglos son iguales en longitud) pero a lo menos tendrá
            0 elementos (el peor de los casos para la altura de recursión, y cómo hay cero elementos, ni siquiera
            se empiezan a apilar cosas en la pila). El lado faltante lo ejecutamos de manera iterativa con el
            flujo natural del while (en el cual no se ocupa apilar nada en la pila). Así vamos a lograr que
            la máxima altura de la pila de recursión sea $\Theta(\log n)$ en el peor caso.

            
            \noindent\begin{minipage}{\linewidth}
            \begin{algorithm}[H]
                \DontPrintSemicolon
                \KwIn{Un arreglo $A$ y dos enteros $p$ y $r$.}
                \KwOut{El mismo arreglo $A$, ordenado.}
                \SetKwProg{Fn}{Algoritmo}{:}{}
                \Fn{$\textsc{OtroQuicksortModificado}(A,p,r)$}{
                    \While{$p < r$}{
                        $q \gets \textsc{Particion}(A,p,r)$\;
                        \eIf{$q-p \leq r-q$}{
                            $\textsc{OtroQuicksortModificado}(A,p,q-1)$\;
                            $p \gets q+1$\;
                        }{
                            $\textsc{OtroQuicksortModificado}(A,q+1,r)$\;
                            $r \gets q-1$\;
                        }
                    }
                }
                \caption{\textsc{OtroQuicksortModificado}}
            \end{algorithm}
            \end{minipage}

            
        \end{enumerate}

    \item Dados dos enteros positivos $n$ y $k$ y un arreglo $A$, describa un
        algoritmo que encuentre los $k$ elementos m\'as peque\~nos de $A$ y los
        devuelva, en orden, en un arreglo de salida $B$.   El algoritmo debe
        tener un tiempo de ejecuci\'on, en el peor caso, de $O(n \log k)$.
        Demuestre que su algoritmo es correcto y la cota de tiempo es la
        deseada.


        \medskip

    \item Entre 1966 y 2004, la UNAM tuvo un sistema de calificaciones con
        letras en lugar de n\'umeros:
        \begin{itemize}
            \item \textsf{MB} - Muy bien, equivalente a 10.
            \item \textsf{B} - Bien, equivalente a 8.
            \item \textsf{S} - Suficiente, equivalente a 6.
            \item \textsf{NA} - No Acreditado, 5 o menos.
            \item \textsf{NP} - No Presentado.
        \end{itemize}
        Exhiba un algoritmo que, dada una lista de $n$ estudientes, la ordene
        por calificaci\'on, i.e., primero a todos los estudiantes con
        calificaci\'on \textsf{MB}, despu\'es a quienes tienen calificaci\'on
        \textsf{B}, etc.   Su algoritmo debe usar tiempo a lo m\'as $O(n)$ y
        espacio $O(1)$.   Adem\'as, su algoritmo debe ser estable (si los
        estudiantes $i$ y $j$ tienen la misma calificaci\'on e $i$ aparec\'ia
        antes que $j$ en la lista original, tambi\'en debe suceder que $i$
        aparezca antes que $j$ en la lista ordenada).

    \item ?`Cu\'antas comparaciones son necesarias y suficientes para ordenar
        una lista con 5 elementos?   Justifique su respuesta.

    \item Dados dos arreglos ordenados $A$ y $B$ de longitud $n$ y $m$,
        respectivamente, exhiba algoritmos de tiempo $O(n+m)$ que calcule:
        \begin{enumerate}
            \item El punto medio de la uni\'on ordenada de ambos arreglos.
            \item La diferencia sim\'etrica de $A$ y $B$ sin repeticiones.
        \end{enumerate}

    \item Considere el siguiente algoritmo.
    
        \begin{minipage}[H]{0.943\textwidth}
        \begin{algorithm}[H]
            \DontPrintSemicolon
            \KwIn{Un arreglo $A$.}
            \SetKwProg{Fn}{Rutina}{:}{}
            \Fn{$\textsc{Algoritmo}(A)$}{
                $i \gets 0$\;
                \While{$i < |A|$}{
                    \If{$i = 0$ \KwOr $A[i-1] \le A[i]$}{
                        $i \gets i+1$\;
                    }
                    \Else{
                        \KwExchange $A[i]$ con $A[i-1]$\;
                        $i \gets i-1$\;
                    }
                }
            }
            \caption{\textsc{Algoritmo}}
        \end{algorithm}
        \end{minipage}
        \begin{enumerate}
            \item Describa en alto nivel qu\'e hace \textsc{Algoritmo} y c\'omo
                lo hace.
            \item Demuestre que \textsc{Algoritmo} es correcto respecto a su
                respuesta del inciso anterior.
            \item Analice la complejidad de \textsc{Algoritmo} en el peor caso.
            \item Analice la complejidad de \textsc{Algoritmo} en el mejor caso.
        \end{enumerate}

    \item Considere el siguiente algoritmo de ordenamiento.   Dado un arreglo
        $A$, para cada posici\'on $i$ en el arreglo, nuestro algoritmo determina
        cu\'al es la posici\'on final $j$ del elemento $A[i]$, esto lo logra
        contando cu\'antos elementos hay a la derecha de la posici\'on $i$ que
        son menores que $A[i]$.   Al intercambiar $A[i]$ con $A[j]$, sabemos que
        el elemento en $A[j]$ ya est\'a en su posici\'on correcta.  Si tras el
        cambio el elemento en $A[i]$ ya est\'a en su posici\'on correcta, el
        \'indice $i$ avanza una posici\'on a la derecha.    Si no, el
        procedimiento se repite hasta que $A[i]$ esté en su posici\'on correcta.
        \begin{enumerate}
            \item Formalice la idea anterior mediante pseudoc\'odigo. Considere
                que puede haber elementos duplicados en el arreglo. El n\'umero
                total de escrituras en memoria debe ser a lo m\'as $|A|-1$.
            \item ?`Cu\'al es el tiempo de ejecuci\'on en el peor caso?
                Justifique su respuesta.
            \item Demuestre que si las entradas del arreglo son los enteros $0,
                \dots, |A|-1$, entonces el algoritmo puede modificarse para
                tener tiempo de ejecuci\'on $O(n)$.
            \item Exhiba un algoritmo que, dado un arreglo de longitud $n$ y
                cuyas entradas son n\'umeros entre $0$ y $n-1$ con exactamente
                un n\'umero faltante y un n\'umero duplicado, encuentre c\'ual
                es el n\'umero faltante y cu\'al el duplicado en tiempo $O(n)$.
        \end{enumerate}

    \item Sea $m$ un entero y $A$ un arreglo con $n$ entradas entre $0$ y $m-1$.
        \begin{enumerate}
            \item Si $m \in O(n)$, ?`es posible ordenar $A$ en tiempo $O(n)$?
            
            \medskip
            \Sol Sí es posible ordenarlo, y usando esta nueva información que sabemos
            podemos usar el algoritmo de CountingSort, que vimos en clase y además tiene la
            característica de que es estable y su tiempo de ejecución es $O(n+k)$ donde $k$ es el rango de los elementos a ordenar, en este caso $k = m$.
            Entonces, si $m \in O(n)$, el tiempo $O(n+k) = O(n+m)$ es absorbido brutalmente por $O(n)$, por lo que podemos decir que sí.

            \item Si $m \in O(n^2)$, ?`es posible ordenar $A$ en tiempo O$(n)$?
            \medskip
            \Sol Aquí si cambia un poco la cosa, que que si intentamos usar Counting, su complejidad
            temporal ahora sería $O(n+m) = O(n+n^2) = O(n^2)$. Pero qué pasa si en este caso usamos
            Raddix sort (también visto en clase) con números escritos en base $n$, entonces el tiempo de ejecución sería $O(d(n+b))$ donde $d$ es
             el número de dígitos y $b$ es la base, en este caso $b = n$. Y como sabemos que $m \in O(n^2)$,
              entonces el número de dígitos $d$ es constante. Y miren a continuación muestro matemáticamente como se representarían :
               si $0 \leq x < n^2$ y $n \geq 2$, escribimos $x = an+c$, donde $a = \lfloor x/n \rfloor$ y $c = x \bmod n$; el cociente y el residuo son sus dos dígitos en base $n$, ambos entre $0$ y $n-1$. Si $m$ supera $n^2$ por un factor constante, pueden necesitarse más dígitos, pero su cantidad sigue siendo constante.
               Entonces el tiempo de ejecución sería $O(d(n+n)) = O(n)$.
            
        \end{enumerate}

    \item Dada una cadena de texto de longitud $n$, sobre un alfabeto de
        tama\~no $k$, exhiba un algoritmo que tras un preprocesamiento de tiempo
        $O(kn)$ pueda responder cualquier consulta respecto a cu\'antas veces
        aparece un s\'imbolo en un rango de \'indices en tiempo $O(1)$, por
        ejemplo, ?`cu\'antas veces aparece la letra 'v' en el rango $[28, 56]$?
        Demuestre que su algoritmo es correcto y cumple con la cota de tiempo
        requerida.
\end{enumerate}
\end{document}
